Brain--computer interfaces (BCIs) must be calibrated for each user, and the
burden is greater for a command with no natural counterpart, such as a
robotic sixth finger. We asked how much calibration stroke patients need and
whether data from other patients reduce it, using a public EEG dataset of
\nSubjects{} patients imagining hand and sixth-finger movements. We first
tested fairly: the published accuracy for hand against sixth finger relies
on spatial filters fitted on all trials before cross-validation, and a
standard pipeline reached \rqOneCspLdaPaperReplicaMean\,\% that way but
\rqOneCspLdaCvLeakfreeMean\,\% when every step was fitted on training trials
only. In pre-registered tests on \rqTwoHTwoAN{} patients, re-centred transfer
from other patients reached \rqTwoRestAhSfRctTslrKOneMean\,\% on three classes
(rest, hand, sixth finger) after one trial of each movement
(\rqTwoMinutesKOne\,min), against \rqTwoRestAhSfTsLrKOneMean\,\% from scratch,
and its area under the learning curve was larger ($p = \rqTwoHTwoAP$). It
then levelled off near \rqTwoRestAhSfRctTslrKAllMean\,\%, while the decoder
from scratch overtook it with the whole pool (\rqTwoRestAhSfTsLrKAllMean\,\%
after \rqTwoMinutesKAll\,min). For adding the sixth finger to a hand decoder,
transfer gave no clear benefit, and the patient's own hand decoder detected
sixth-finger imagery at least as well, so we do not claim a distinct
sixth-finger command. In exploratory analyses, a sixth-finger-against-rest
switch reached \exFourTsLrCentralMean\,\% but also responded to imagined hand
movement, and decoders lost accuracy on later days, most of which five new
trials of each movement recovered. Transfer gives a fast start, not a better
decoder.
